% TOP_PROBLEM: {"id": 1407, "title": "Turán's Brick Factory Problem", "category_id": 3, "category": {"id": 3, "name": "graph_theory", "display_name": "Graph Theory", "description": "Problems involving graphs, networks, and their properties.", "slug": "graph-theory", "order_index": 3, "created_at": "2026-07-31T15:26:25.671Z"}, "status": "open"}
% FIRST_POSTED: null
% ENTRY_KIND: statement_only
% Imported from: unsolvedmath_additional_500.tex; UnsolvedMath ID 1407
% !TeX program = lualatex
% Compile twice: lualatex 1407.tex
% Self-contained: no external bibliography, images, or \input files.
% Numeric section labels and visible numbers are original UnsolvedMath IDs.
\documentclass[11pt,a4paper]{article}
\usepackage[margin=24mm,headheight=14pt]{geometry}
\usepackage{fontspec}
\setmainfont{Latin Modern Roman}
\setsansfont{Latin Modern Sans}
\setmonofont{Latin Modern Mono}
\usepackage{amsmath,amssymb,mathtools}
\usepackage{unicode-math}
\setmathfont{Latin Modern Math}
\usepackage{microtype}
\usepackage{enumitem,xurl,needspace}
\usepackage[dvipsnames]{xcolor}
\usepackage{hyperref}
\hypersetup{unicode=true,colorlinks=true,linkcolor=MidnightBlue,urlcolor=MidnightBlue,
 pdftitle={UnsolvedMath Problem 1407},pdfauthor={Source-annotated compilation},bookmarksnumbered=true}
\usepackage{bookmark,tocloft,titlesec,fancyhdr}
\setlength{\cftsecnumwidth}{4.2em}
\cftsetpnumwidth{2.5em}
\cftsetrmarg{4em}
\titleformat{\section}{\Large\bfseries\raggedright}{\thesection}{0.7em}{}
\pagestyle{fancy}\fancyhf{}
\fancyhead[L]{\small\sffamily Additional UnsolvedMath statements}
\fancyhead[R]{\small Original catalogue IDs}
\fancyfoot[C]{\thepage}
\setlength{\parindent}{0pt}
\setlength{\parskip}{5pt plus 1pt}
\setlength{\emergencystretch}{3em}
\setcounter{tocdepth}{1}\setcounter{secnumdepth}{1}
\setlist[itemize]{leftmargin=3em,itemsep=4pt,topsep=4pt,parsep=0pt}
\allowdisplaybreaks
\newcommand{\N}{\mathbb{N}}
\newcommand{\Z}{\mathbb{Z}}
\newcommand{\Q}{\mathbb{Q}}
\newcommand{\R}{\mathbb{R}}
\newcommand{\C}{\mathbb{C}}
\newcommand{\F}{\mathbb{F}}
\newcommand{\A}{\mathbb{A}}
\newcommand{\PP}{\mathbb{P}}
\newcommand{\E}{\mathbb{E}}
\newcommand{\eps}{\varepsilon}
\newcommand{\dd}{\,\mathrm d}
\newcommand{\sref}[2]{\hyperlink{source:#1:#2}{[S#2]}}
\newif\ifshowcatalogue
\showcataloguetrue
\newcommand{\cataloguescope}[1]{\ifshowcatalogue
 \paragraph{Statement recorded in the supplied source.}
 #1\fi}
\begin{document}
% UnsolvedMath ID 1407; source catalogue code GRAPH-020

\setcounter{section}{1406}

\section{The Crossing Number of Complete Bipartite Graphs}\label{1407}

{\small\sffamily UnsolvedMath ID 1407 · Graph Theory}

\subsection{Definitions and mathematical statement}

\paragraph{Shared notation.}Unless a section says otherwise, $\N=\{1,2,\ldots\}$,
$\Z$ is the ring of integers, $\Q,\R,\C$ are the rational, real, and complex
fields, $[N]=\{1,\ldots,\lfloor N\rfloor\}$, and $\log$ is the natural
logarithm. For real functions of a parameter tending to infinity,
$f\ll g$ and $f=O(g)$ mean $|f|\leq Cg$ eventually for some fixed $C$;
$f\gg g$ reverses this comparison; $f=o(g)$ means $f/g\to0$;
$f\sim g$ means $f/g\to1$; and $f\asymp g$ means both order bounds.
A subscript on $O$ or $\ll$ specifies allowed constant dependence.
The expression $n^{o(1)}$ in an upper bound means growth smaller than
$n^\epsilon$ for each fixed $\epsilon>0$. Local definitions govern any
exception, finite-field convention, or use of the same symbol for another
object. An asymptotic claim is not automatically an effective algorithm.



The graph $K_{m,n}$ has disjoint vertex classes of sizes $m,n$ and every possible edge between the classes, with no edges within a class. Crossing number permits arbitrary topological drawings. The familiar proposed value is
\[Z(m,n)=\lfloor m/2\rfloor\lfloor(m-1)/2\rfloor\lfloor n/2\rfloor\lfloor(n-1)/2\rfloor.\]
The source asks for the exact minimum, not merely the value of one explicit drawing. \sref{1407}{1} \sref{1407}{2}

\paragraph{Statement recorded in the supplied source.}
What is the minimum crossing number of the complete bipartite graph $K_{m,n}$? \sref{1407}{1}

\paragraph{Residual task reported by the archive.}

The embedded review (2026-08-17) records: Prove Zarankiewicz's lower bound for all m,n or find a drawing with fewer than Z(m,n) crossings. \sref{1407}{1}

\subsection{Short English statement}

What is the fewest number of crossings needed to draw all connections between two groups of prescribed sizes?

\subsection{Sources}

{\small

\begin{itemize}

\item[\hypertarget{source:1407:1}{S1}] ULAM.AI, \emph{UnsolvedMath}, supplied \texttt{problems.json}, record 1407 (GRAPH-020), statement and background. The embedded literature-review date is 2026-08-17; that is the archive’s review date, not a fresh certification. Dataset landing page: \url{https://huggingface.co/datasets/ulamai/UnsolvedMath}.

\item[\hypertarget{source:1407:2}{S2}] J. Balogh, B. Lidický, S. Norin, F. Pfender, G. Salazar, and S. Spiro, Crossing numbers of complete bipartite graphs, Procedia Computer Science 223 (2023), 78–87. \url{https://doi.org/10.1016/j.procs.2023.08.216}. Bibliographic information imported from the supplied record.

\item[\hypertarget{source:1407:3}{S3}] E. de Klerk and D. V. Pasechnik, Improved lower bounds for the 2-page crossing numbers of K\_{m,n} and K\_n via semidefinite programming, arXiv:1110.4824. \url{https://arxiv.org/abs/1110.4824}. Bibliographic information imported from the supplied record.

\end{itemize}

}

\label{end:1407}
\end{document}
